% Version 3.9: public import, local request gossip and regional first-DATA reports.
% Replaces sec:requests up to, but not including, Faithful publication.
\subsection{Requests, input reports, and retained data}
\label{sec:requests}
A request is \defterm{admissible} when it passes the application's admission
checks. Its logical identity distinguishes a retransmission from a new
client command. Let $\operatorname{enc}(x)$ be its canonical authenticated
representation: transport wrappers and alternate encodings of the same
authentication proof do not create a new request. Each node retains
requests in first-receipt FIFO order; repeated delivery does not change
that order. Different nodes need not have the same order.

A request is finally removed only when its identity occurs in $G$.
Regional commitment, regular handoff and discarded speculative suffixes
do not remove it. The first-DATA primary rule below replaces the original
oldest-pending selection at that position. After the primary, and at
later positions, a correct leader fills the remaining capacity by its
ordinary local FIFO, skipping identities already in the exact parent
ancestry or batch and rechecking application validity after earlier batch
entries. This is a change to correct-leader batching, not a stronger
input-delivery or common-order assumption.

\paragraph{Request relay within an installed snapshot.}
On first complete receipt of an authenticated admissible request $x$, a
correct node persists its pending entry. For each authenticated installed
snapshot of its region in which it is a member, it dispatches the full
request to that snapshot's members, once per request and snapshot.
Learning or installing such a snapshot also schedules the same individual
relays for retained requests. Receivers apply the same first-receipt rule.
An authenticated relay is not a certificate that $x$ is admissible after
every future parent; that check remains local to batch construction.

The request store and pending relay actions have a persistent owner in
Waiting, Ready, Active, Frozen, Completed and Terminated. Closed and
\texttt{RequestClose} do not delete them or block request transport.
This action signs no BFT vote, changes no lock, and starts no timer or
entry. It adds no Ready gate and waits for no whole-pool transfer.
For one installed snapshot, a correct member's dispatched full request
reaches and is processed by every correct member within its existing
$\Delta_i$ bound. In particular, immediate first-receipt dispatch gives
this bound from that complete receipt handler. This is a bound on the
specified within-snapshot messages, not on discovering a new roster,
dispatching an arbitrary backlog, or transferring between disjoint
rosters. The latter awareness obligation remains separate below.

% Version 3.9: nonblocking request import from retained public DATA bytes.
% Insert after the installed-snapshot relay and before INPUT_REPORT.
\paragraph{Request import from retained public DATA.}
\label{sec:public-request-import}
Each correct node's persistent request owner also runs the local import
loop below. Its input is the actual bytes already retained in
\tagg{PublicStore}, not a request query to an old roster. Its output is
a store of unvalidated request candidates and, after the ordinary
request checks, entries in the existing pending FIFO. Import neither
validates its containing DATA nor changes $G$, a boundary, a lock, a
vote, a report or an entry gate. In particular, it does not reopen a
closed session.

The owner retains an append-only source-work journal $Q^{\rm src}_v$,
an append-only candidate store $U_v$, and an append-only checking journal
$Q^{\rm chk}_v$. Each journal has a durable completion cursor and, for
its current finite job, a saved local subposition. Candidate entries
retain their exact bytes, source-object occurrence paths and local
arrival numbers. Checking outcomes and pending-store updates are
durable. A completed job is marked complete; its evidence is not
deleted. Existing pending entries retain their first-receipt order.

When complete raw source content is retained, append its actual retained
root and dictionary cut to $Q^{\rm src}_v$ before releasing the source
input or its enqueue obligation. Do the same when an exact public
reference target is subsequently acquired. On first enabling this rule,
enqueue one fixed finite cut of the already retained roots; concurrent
acquisitions append separate jobs. This is a local scan of held bytes,
not a new hydration or Ready condition. The new journal cursor is
independent of source availability, P6 and regional entry cursors.
Retaining the descriptor does not run its scan in a timed handler.

A source job follows the original typed fields and list positions of
\cref{sec:logical-carrier}, including DATA in proposals, ancestry and
boundary evidence. In the existing block tuple
\[
 B=(\beta,r,k,h,\mathit{type},\mathit{batch},
          \Hash(B_{\rm parent}),\rho,\xi),
\]
it extracts the actual entries of $\mathit{batch}$ only when the local
structure identifies a DATA body and supplies those entry bytes.
DRAIN has no batch. Arbitrary witness strings, request hashes and a
sender's object index are not DATA batch entries. Decode independently
represented typed children even if a sibling or enclosing proof fails.
Malformed bytes without that logical structure yield no invented entry.

The job traverses supplied inline objects and exact public targets held
at its fixed dictionary cut, visiting each actual object once. It saves
an unresolved edge without waiting for or fetching its target; a later
acquisition supplies another job. A cycle cannot cause repeated visits.
Thus the job has a finite fixed input. It does not fetch AUTH, native
state, source history or optional proof targets. Outer chain, version,
base and interval hints, P6 rejection and branch closure do not filter
this scan. An actual DATA entry is still only a candidate, including
when the containing DATA claims the wrong signed domain.

Append each new exact candidate byte string to $U_v$ before advancing
past its batch position. Before successful request validation,
deduplicate only identical bytes, with byte equality or a binding
content commitment, and preserve their source occurrences. A claimed
client identifier creates no duplicate marker. Invalid encodings and
failed authentication cannot reserve an identifier or suppress a later
valid request. $U_v$ is explicitly unvalidated: its entries are not
eligible reports, authenticated relays or executable requests.

\paragraph{Checking candidates without a moving target.}
When the ordinary application interface has completely prepared an exact
parent $P$, enqueue a checking job for the finite current cut of $U_v$.
When a new candidate is retained and such a parent is already prepared,
enqueue a check of that candidate at that parent. Equal
candidate-byte/parent checks need not be repeated. Each job retains the
already available read context needed by the ordinary finite request
check, including its exact native state, until the check finishes.
A pointer to a subsequently mutated engine state is not such a context.
Any required snapshot construction must actually finish before this
notification; neither its cost nor instantaneous copying is inferred
from the import rule. Ready does not wait for this construction.
Changing the active parent does not retarget or cancel an old job.
These jobs acquire no absent validation input and cause no state
reconstruction; preparing a parent keeps its original prerequisites.

At that fixed $P$, run the ordinary chain/domain, client authentication,
canonicalization and application-admission checks. A failed check is an
outcome only for those bytes at that parent, not a permanent rejection
of an asserted identity or a substitute for checks at later parents.
A success establishes $\operatorname{enc}(x)$ and the request's actual
logical identity under the original request interface. Only now merge
retransmissions by that validated identity; an asserted equal identifier
alone is insufficient. If its identity is already in $G$, record the
outcome without adding an active pending entry. Otherwise persist it in
the ordinary pending store, preserving an existing entry's FIFO position
or assigning the next first-receipt position to a new entry, and invoke
the same installed-snapshot relay rule. A successful check at an old
parent is not admission at a newer one.

The loop continuously alternates finite local steps on the oldest
unfinished source job and the oldest unfinished checking job, skipping
an empty journal. Within each job it follows the saved field or
candidate order. New jobs go to the tail and cannot overtake a fixed
job. The node continues these steps until each queued finite job is
complete; session changes and repeated new arrivals cannot suspend it
indefinitely. This is an explicit eventual local-work duty, with no
wall-clock deadline, throughput bound or finite-storage claim.
Neither queue waits on a fetch or on the other queue's completion.
Ordinary messages, timers, P6, Ready and report emission do not wait for
the import loop, its backlog, or an import watermark. Their original
processing bounds are not bounds on this deferred work.

The journals, bytes, read contexts and pending enqueue obligations belong
to the node, independently of a producing engine or import worker.
Waiting, Frozen, Completed, Terminated, \texttt{RequestClose}, Closed,
checkpoint compaction and advancement of $G$ cannot drop them. A worker
exit transfers them durably and a restart resumes their exact cursors.
The existing public-source retention rules continue to govern their
underlying bytes. Inclusion in $G$ removes an active pending request;
it does not erase raw evidence or turn a rejected record into accepted
history.

\paragraph{Interface to reports and FIFO filling.}
The set $E_v(P)$ below is still the exact eligible projection of the
authenticated pending store. An unvalidated entry of $U_v$ is not yet a
pending request; deferred checking is not permission to omit an admitted
pending request from $E_v(P)$. A successful import enters the ordinary
request-processing interface, which completes the required exact-parent
admission and ancestry filtering before exposing the updated projection.
At report emission and FIFO filling, use the completely validated actual
parent, its native state and complete ancestry, including inherited
global history and speculative ancestors. A previous import check or
public occurrence replaces none of those checks. Earlier entries in a
batch still affect admission of the remainder. An emitted report remains
immutable; an import completed later is available to later report
contexts and ordinary FIFO filling.

\begin{proposition}[Import after a first public carrier]
\label{prop:v39-public-request-import}
Suppose an original complete public result $R_0$ contains the complete
bytes of a request $x$ in an actual DATA batch in its old logical fields
or unchanged public source denotations. If a correct node follows and
retains those actual bytes and continues the local loop, it eventually
retains $x$ in $U_v$, even if the record is rejected or its branch is
closed. If, after that retention, an exact parent at which $x$ passes
the ordinary request checks is completely prepared, or such a parent
is already prepared then, the request eventually enters the pending
store unless its identity has already entered $G$.
\end{proposition}
\begin{proof}
The original public fields determine the DATA body and its actual batch
entries. Source retention records a job for each acquired root or exact
target. Each fixed job has only finitely many predecessors, a finite
input and finite local traversal; tail arrivals cannot overtake it.
The prescribed loop therefore reaches the actual entry and persists its
bytes. Failed hints, branch closure and unrelated witness failures are
not scan gates. An unresolved edge alone supplies no bytes, but the
hypothesis supplies acquisition of its actual target and hence its job.
Exact-byte deduplication only finds a copy already retained; an invalid
claimed identifier cannot suppress the entry.

The prepared-parent event or candidate arrival likewise records a fixed
finite checking job. Its retained local context needs no fetch. FIFO
service eventually completes the ordinary successful check of $x$ and
updates the pending store. A later parent cannot cancel that work.
Thereafter the ordinary pending rule retains the request until its
identity enters $G$, while every report and batch independently checks
eligibility at its own exact parent.
\end{proof}

This result gives no bound from public visibility to pending admission.
In particular, valid candidate bytes cannot be omitted forever merely
by leaving their finite checking job unprocessed, but import need not
finish before any particular report. The proposition neither forms
the first $x$-bearing $R_0$ nor ensures that every future roster knows
$x$ before its first DATA. It adds no request-publication service,
cross-snapshot message guarantee or historical quorum-availability
premise. Accepted-report awareness remains a separate condition of
\cref{prop:v38-request-inclusion}.


\subsubsection{The first-DATA report}
\label{sec:input-report}
Order requests by
\[
 \operatorname{key}(x)=
   \bigl(|\operatorname{enc}(x)|,\Hash(\operatorname{enc}(x))\bigr),
\]
lexicographically, with $\tagg{EMPTY}$ after every request key. Hash binding
identifies the actual canonical body. Every fixed request has finitely
many predecessors in this order, since they are byte strings of bounded
length. This order provides no practical latency bound.

Unvalidated public candidates are kept separately as in
\cref{sec:public-request-import}. For a completely validated parent $P$,
let $E_v(P)$ contain the pending
requests whose full bytes node $v$ holds, that are admissible immediately
after $P$, fit alone in a DATA batch, and whose identities are absent from
the complete ancestry of $P$. The ancestry includes the inherited prefix
and uncommitted speculative blocks. Maintain a key-ordered index of
retained requests and its eligible projection for the exact parent.
Parent and request processing completes the required admission and
ancestry filtering before exposing this projection to the report action.
The index is never reused for another parent hash or native state.

This selection uses the same abstract batch-processing scope as the
original pending-admissible batch builder. It introduces no
\texttt{AcquirePending} query or acquisition wait. The complete-handler
timing below does not bound an extra deferred pool scan, native import
or background eligibility rebuild. If an implementation leaves such
work outside completed message/batch processing, its old timing bound
has not thereby been proved. An index alone is not a measured CPU or
bandwidth guarantee.

\paragraph{Context and emission.}
For $V=(\beta,r,g,t)$, write $P=A_{\beta,r,g,t}$. A report signs the chain,
protocol domain, type $\tagg{INPUT\_REPORT}$, session, assigned region,
configuration, $\tagg{AUTH}_{r,g}$ reference, $(g,t,j=0)$, the generation's
opening identity, and $P$'s hash, DATA height and physical position,
together with one request key or $\tagg{EMPTY}$. Denote this full context
by $\mathit{RC}$. Its boundary and authority are authenticated by the
original admission rules; a claimed pair alone is insufficient.
Alternate QC encodings do not change $\mathit{RC}$.

A correct member emits its report in the ordinary-term entry handler
for $t>0$, after validating the preceding terminal QC and exact parent.
For $t=0$, it emits on complete validation of the NEWVIEW in the entered
generation, after the exact opening is known. Require the original
fresh-action gate, an entered current unfrozen generation, and a
DATA-producing term. Persist one report for $\mathit{RC}$ before sending:
choose $\min E_v(P)$, or $\tagg{EMPTY}$ if this set is empty, and send the
signed report and the chosen request's full authenticated bytes to
$L(V)$. Resends repeat that report. Later arrivals do not amend it.
Historical validation never creates a missing fresh report.

The original immediate NEWVIEW forwarding and ANCHOR vote remain
independent. They wait for neither a report nor a report certificate.
If both messages share a packet, parse them independently: a bad report
does not suppress a valid ANCHOR vote or original Quit evidence.
Reports create no QC rank, lock, timer, commitment, entry permission or
new Quit proof. Fresh report signing obeys retirement; retaining and
relaying existing reports do not reopen its signing gate.

\paragraph{Correct-leader collection and batching.}
For the exact $\mathit{RC}$, a correct leader retains distinct-member
reports only after checking their signatures, snapshot and context.
For each nonempty retained report it also checks the full body, key,
admission after $P$, individual fit and ancestry exclusion. Ignore
incomplete or invalid contributions; count a member at most once.
Once these reports have weight $>F$ and the original parent/ANCHOR QC
is available, construct the first DATA as follows. If any reported key
is nonempty, execute the minimum reported request first as the
\defterm{primary}, then fill by ordinary FIFO. Otherwise use ordinary
valid batching, possibly empty. New requests may therefore be included
even when all saved reports are EMPTY. A singleton primary fits by its
eligibility test. Correct reports alone have sufficient weight; the
leader never waits for Byzantine replies or an identifiable event
called ``all correct replies received.''

Attach the signed report keys, contexts and member identities as a
separate regional validation sidecar. The selected body is already in
the ordinary DATA batch. Other validators need not obtain nonselected
request bodies. The original signed DATA, PROPOSE, votes, QCs and FINAL
objects, and their hash encodings, are unchanged.

\paragraph{Full regional validation.}
For every DATA proposal at $j=0$, require the original full validation
and a report set of distinct snapshot members with weight $>F$ for its
exact historical $\mathit{RC}$. If any key is nonempty, check that its
minimum is the batch's first request, with matching authenticated bytes,
admission immediately after $P$, individual fit and no identity in the
parent ancestry. Check the entire batch and exclude another occurrence
of the primary identity. For all-EMPTY reports apply ordinary batch
validity; an empty batch is allowed, not required.

Validators do not compare the witness with their own queues, require
their own report, or apply current-time/current-$G$ tests to a historical
report. Correct reporter selection supplies the local-minimum fact;
the common predicate checks the supplied signed minimum and selected
body. A fictitious selected key cannot pass that body check.
For later DATA positions and DRAIN, validate the required first-DATA
sidecars in their own-session ancestry; DRAIN creates no new primary.
Apply this rule in Vote, recovery histories, own-session receipt
processing and Completion. Earlier-term current-generation proposals
retain the original Vote eligibility after local advancement. Old and
future generations retain their original archival/buffering behavior.

Evidence-first processing still inspects original signed material even
when a report wrapper is invalid or incomplete. Report validation must
not hide an existing admissible DATA/DRAIN or NEWVIEW equivocation.
Different sidecar signer sets or encodings for the same proposal are
witness variants: proposal identity, deduplication and timer identity
remain $(\beta,r,g,t,j,\Hash(B))$.

\subsubsection{Regional evidence and the original public result}
\label{sec:input-report-carriage}
Write $R_0$ for the original complete regional result of
\cref{sec:delivery}, with every original DATA, FINAL, proposal, QC,
ancestry, boundary, authority, native-validation and public-source object
required by that interface. It is not merely the public kernel.
An augmented regional bundle is
\[
 R_+=(R_0,W_R),
\]
where $W_R$ supplies all report sidecars needed for augmented own-session
validation. Correct regional senders forward and retain complete $R_+$;
STATUS/NEWVIEW histories, receipt relays, Completion and retirement
ownership transfers preserve the applicable $W_R$. References must be
resolvable within the existing complete-message bound. A receipt's
original public representation alone need not provide augmented
own-session Completion inputs.

The Publish Layer's input remains exactly $R_0$. Erasure removes only
the new separately framed sidecars; it preserves all old required
objects, signatures and established public denotations. For the same
DATA and original witness, the public record is the same old-format
record, with zero additional REPORT bytes. The primary occupies ordinary
application capacity. Original public size/load limits and the original
$T_{\rm pub}$ starting event remain unchanged; no service bound applies
to bare requests or reports.

The honest-FINAL argument now traces an actual correct signature to a
timer-origin Vote that checked both original validity and the report
predicate. It therefore proves existence of a genuine historical
report witness without making it a public fetch dependency. FINAL still
binds no generation or term and does not authenticate a publisher's
chosen alternate witness. Cross-session inherited anchors retain the
historical kernel rule. All independent global extension, branch,
reference and Force-cutoff checks remain necessary.

Keep block data and receipts needed for recovery and checkpoint
verification. Before finalized Espresso inclusion, keep complete $R_0$
results available to the Publish Layer and direct readers. After
finalization, other regions may obtain them from authenticated Espresso
history. Regional sidecars remain with their required data/evidence
owners; public erasure is not permission to release their last required
copy. Global promotion alone authorizes no archival deletion.

\subsubsection{Preserved timing and conditional request progress}
\label{sec:input-report-progress}
\begin{proposition}[Local report timing]
\label{prop:v38-report-timing}
Under the original complete-handler scope, entry preparation and active
unfrozen/unclosed interval of the correct-leader progress argument, the
new rule meets the original fresh-position BLAME deadlines. Ordinary
first DATA is voted by all correct members within $4\Delta$ of first
correct term entry; initial/recovery term zero retains $6\Delta$.
\end{proposition}
\begin{proof}
Let $e$ be first correct entry. In an ordinary term, the preceding
terminal proposal's complete forwarding and correct votes disseminate
its QC to all correct members by $e+2\Delta$, independently of the
successor's proposal. Every correct report is sent at that local entry
and reaches a correct leader by $e+3\Delta$. Correct weight exceeds $F$;
the complete first proposal is voted everywhere by $e+4\Delta$.

In initial/recovery term zero, the original prepared entry skew is
$\Delta$. A correct initializer sends its NEWVIEW by $e+3\Delta$;
complete processing everywhere by $e+4\Delta$ sends the independent
ANCHOR votes and reports. Both correct-only inputs reach the leader by
$e+5\Delta$, and first DATA is voted by $e+6\Delta$. DRAIN adds no report.
Every later position costs at most $2\Delta$. The $p$th fresh vote is
therefore no later than $e+(2p+4)\Delta$, and hence no later than each
member's original deadline $e_v+(2p+4)\Delta$. At equality messages
precede progress callbacks. Arbitrary authority/native preparation
waits retain their original eventual scope; no new bound is asserted.
\end{proof}

Complete forwarding carries the same common report predicate to every
correct recipient, without a private-queue veto. The original
protecting-QC, lock and timer arguments therefore still apply to an
actual augmented-valid Vote. The old last-check bound for a faulty term
also applies: report collection never suppresses progress checks, and a
correct terminal vote still forwards every dependency needed for other
correct votes. A correct initializer still proposes by $e+5\Delta$,
retains its protected commit window and forms a complete receipt by
$e+9\Delta$. Thus the existing prepared regional DATA-response bounds
survive this rule. This reasoning replaces the sharper old ordinary
handoff latency; it does not retain a $2\Delta$ first-vote claim.

\begin{proposition}[Request inclusion from accepted report awareness]
\label{prop:v38-request-inclusion}
Fix a request $x$ that fits alone and remains admissible after the
relevant accepted parent states until inclusion. Suppose infinitely
many accepted term-first DATA have a genuine report witness containing
a correct signer whose eligible retained set contained $x$ at its
report event. Then $x$ enters $G$. The implication permits arbitrary
changes of roster between sessions.
\end{proposition}
\begin{proof}
Suppose $x$ never enters $G$. At any such accepted first DATA, its exact
parent ancestry cannot contain $x$: accepting that ancestry would
already include $x$. The aware correct reporter chooses $y$ with
$\operatorname{key}(y)\le\operatorname{key}(x)$; the certificate minimum
is an actual executed request $z$ with
$\operatorname{key}(z)\le\operatorname{key}(y)$.
Its identity is absent from the exact parent ancestry. Upon acceptance
that ancestry is the global prefix preceding the new DATA, including
ancestors accepted in the same EXPORT, so $z$ spends a previously
unaccepted predecessor identity.

There are finitely many canonical byte strings with key at most
$\operatorname{key}(x)$. Let $\Phi_x(G)$ count their distinct logical
identities absent from $G$. It is finite, never increases as $G$
extends, and strictly decreases at each counted accepted first DATA
while $x$ is absent. Infinitely many such decreases are impossible.
Discarded speculative commitments do not spend this potential; accepted
predecessors cannot be spent twice because every later accepted parent
contains them. Recovery and Force never roll back $G$.
\end{proof}

Each DATA-producing term has at most $K$ positions. Every genuine DATA
history contains its term's first DATA at most $K-1$ DATA steps earlier;
contiguous acceptance cannot skip that ancestor. Thus unbounded accepted
DATA contains infinitely many distinct accepted term-first blocks,
even when recovery changes their historical proposal contexts. The
same counting applies to a region's repeated accepted contributions.
It remains necessary to prove awareness for sufficiently many of these
genuine witnesses, not for a publisher's claimed term annotation.

Immutable reports are compatible with this implication. A late request
can leave an old EMPTY certificate usable for its exact term, but that
context has only finitely many DATA positions. A finite prefix of
execution events contains only finitely many pre-dissemination reports;
delayed publication cannot turn that stock into infinitely many
term-first blocks. A changed parent, generation, term or branch/session
requires the corresponding new report context.

This does not yet derive full request inclusion from the frozen
changing-membership ingress prose. The within-snapshot relay above
establishes awareness after an actual member's complete receipt, and
retention supports its later reports. It does not prove that disjoint
future rosters receive $x$ before issuing their first reports. Infinitely
many such rosters could each contribute an EMPTY-report block before
their own later delivery event. Per-node eventual delivery alone does
not rule out that order. The candidate leaves this node-level
cross-roster obligation explicit, without a new ownership gate,
whole-pool service, public request publication or stable-membership
restriction.
